\chapter{Procedure} \label{ch:Procedure}

\section{Measurement Method} \label{sec:MeasurementMethod}

The experimental investigation employs a Geiger-Müller counter system configured for counting ionizing radiation events from a radioactive source. Each decay event that triggers the detector produces a distinct electrical pulse, which is registered by the connected counting electronics. The fundamental measurement principle involves recording the number of pulses accumulated during predefined time intervals, thereby converting the inherently continuous decay process into discrete count data suitable for statistical analysis. Multiple series of measurements are conducted under varying conditions to characterize different aspects of the decay statistics. The primary measurement series consists of repeated consecutive countings with fixed duration to establish the count distribution and verify its conformity to Poisson statistics. Additional measurement sequences explore the temporal evolution of count rates over extended periods to assess stability and drift effects. To quantify background contributions, dedicated background runs are performed without the radioactive source present, capturing ambient radiation levels from cosmic rays and environmental radioactivity. Geometric dependencies are investigated by systematically varying the source-to-detector distance while maintaining all other parameters constant, allowing verification of the inverse-square law behavior. Throughout all measurements, careful attention is paid to avoiding parallax errors, ensuring reproducible positioning, and minimizing external disturbances that could compromise measurement integrity. The collected raw data consist of count numbers paired with their corresponding accumulation times, forming the basis for subsequent uncertainty analysis and statistical testing.

\section{Experimental Setup} \label{sec:ExperimentalSetup}

The measurement apparatus comprises several interconnected components assembled to facilitate accurate detection and registration of ionizing radiation. The central detection element is a Geiger-Müller tube housed within a protective shielding enclosure, designed to maximize sensitivity to beta and gamma radiation while minimizing interference from charged particles. Connected to the detector head is a high-voltage power supply that establishes the necessary electric field within the gas-filled tube for avalanche multiplication of ionization electrons. The pulse output from the detector feeds into a scaler-timer unit capable of accumulating counted events and displaying elapsed measurement durations with digital precision. All electronic components are interconnected using standardized coaxial cables to minimize signal loss and electromagnetic interference. The radioactive source, securely mounted on an adjustable stand, can be positioned at precisely controlled distances from the detector window through a calibrated translation stage. Lead shielding blocks are available for background measurements and for attenuating stray radiation during alignment procedures. The entire assembly rests on a vibration-damped optical table to prevent mechanical perturbations that might affect geometric reproducibility. Power connections are organized to allow emergency shutdown capability while maintaining stable operating voltages throughout extended measurement sessions.

% \subsubsection*{Sketch of Experimental Setup} \label{sec:Sketch}
\img[0.7]{Versuchsaufbau}{Experimental setup}

\section{Spectrum for each single element}
For this protocoll eight different elements have been evaluated. All eight elements and ther messured spectrum can be found in in figure \ref{fig:Bild_alle}
\img[1]{Bild_alle}{All elements and ther spectrums.}

The messured elements are (Fr, Mo, Ag, Cu, Ni, Ti, Zn, Zr) -- with their atomic number $Z = (26, 42, 47, 29, 28, 22, 30, 40)$.
